% ====================================================================
%  SL-SPV — MPhil thesis master file
% --------------------------------------------------------------------
%  Compile with:
%      cd voxceleb_trainer
%      lualatex thesis_master
%      bibtex  thesis_master
%      lualatex thesis_master
%      lualatex thesis_master
%  Or use latexmk:
%      latexmk -xelatex thesis_master.tex
%
%  Uses DejaVu Serif (system font) to match the look of the audit PDF.
%  Switch to your institution's required class file at submission time;
%  the chapter sources require only \chapter / \section / \cite, all
%  of which are provided by `report` and any thesis-class derivative.
% ====================================================================

\documentclass[12pt,a4paper,openany]{report}

% Engine-adaptive font setup. The May 2026 build used xelatex +
% fontspec (DejaVu Serif); this machine's TeX Live has since lost
% xetex and its OpenType font support, so the pdfLaTeX branch (Latin
% Modern) is the working path. The fontspec branch is kept for
% submission-time builds on a full TeX Live.
\usepackage{iftex}
\ifPDFTeX
  \usepackage[T1]{fontenc}
  \usepackage[utf8]{inputenc}
  \usepackage{lmodern}
\else
  \usepackage{fontspec}
  \setmainfont{DejaVu Serif}
  \setsansfont{DejaVu Sans}
  \setmonofont{DejaVu Sans Mono}
\fi

\usepackage{geometry}
\geometry{margin=1in}

\usepackage[bookmarks=true,colorlinks=false]{hyperref}
\usepackage{booktabs}
\usepackage{tabularx}
\usepackage{amsmath,amssymb}
\usepackage{microtype}
\usepackage{enumitem}
\usepackage{parskip}

\title{Low-Resource Speaker Verification for Sinhala and Tamil\\
\large An open-source toolkit, a reproducible benchmark, and a
single-corpus empirical study}
\author{[Your Name]}
\date{\today}

\begin{document}

\maketitle

\begin{abstract}
\noindent
Speaker verification (SV) systems have converged on architectures and
training recipes optimised for high-resource English benchmarks;
Sinhala and Tamil, spoken by over 22 million Sri Lankans, have no
published SV baselines. This thesis contributes (i) a reproducible
527-speaker Sinhala/Tamil benchmark built from public read-speech
corpora (OpenSLR SLR52/SLR65), with an in-the-wild extension via our
group's 280-speaker SLCeleb corpus; (ii) the first zero-shot SV
baseline table for Sinhala: four modern VoxCeleb-trained
architectures degrade $2$--$8\times$ in EER under the language shift,
with badly calibrated scores throughout; (iii) a training-free
adaptation study showing AS-Norm alone recovers up to 39\% relative EER
(best training-free system: ReDimNet-B6 + AS-Norm at 2.07\%/0.90\%
EER si/ta), that PLDA's value is condition-dependent, and that
cross-language calibrator transfer inflates $C_{\mathrm{llr}}$ up to
$5.7\times$ — per-language calibration is mandatory; and (iv) a
four-policy fine-tuning comparison of the WavLM-base+ foundation
model in which full fine-tuning (1.69\%/1.41\% EER over three seeds)
decisively beats LoRA on this benchmark, contrary to recent PEFT
findings, with the closed-set caveat analysed. The underlying
toolkit extends the open-source \texttt{voxceleb\_trainer} codebase
with eleven feature modules and twenty-six bug fixes, and all trial
lists, configurations, and result artifacts are released for
downstream Sri Lankan SV research.
\end{abstract}

\tableofcontents

% ====================================================================
%  Chapter 1 — Introduction
% --------------------------------------------------------------------
%  Template chapter. Replace every bracketed [LIKE_THIS] placeholder
%  with the actual value before submission. The chapter sections follow
%  the standard MPhil/journal-paper introduction structure:
%    1.1 Motivation
%    1.2 Problem statement
%    1.3 Research questions
%    1.4 Contributions
%    1.5 Thesis structure
%  Citations resolve against the .bib block in RUN_GUIDE.md §16.
% ====================================================================

\chapter{Introduction}
\label{ch:introduction}

\section{Motivation}
\label{sec:motivation}

Speaker verification (SV) — the task of deciding whether two audio
recordings come from the same speaker — is a foundational biometric
technology underpinning voice authentication, forensic audio analysis,
and conversational-AI personalisation. The dominant publicly available
benchmarks (VoxCeleb1 \cite{nagrani2017voxceleb}, VoxCeleb2, NIST SRE
series) are constructed almost exclusively from English speech. As a
result, the state-of-the-art has converged on architectures and
training recipes that, while highly effective on English, are not
guaranteed to generalise to low-resource languages.

Sri Lanka presents a particularly underserved case. The country's two
primary languages —~Sinhala (an Indo-Aryan language with roughly
17 million speakers) and Tamil (a Dravidian language with roughly
5 million Sri Lankan speakers) —~have, prior to this work and our
group's SLCeleb corpus \cite{slceleb2025}, no publicly catalogued SV
benchmarks and no published SV baselines. Existing
multilingual self-supervised speech encoders cover them only as part
of broader low-resource batches \cite{chen2022wavlm,
baevski2020wav2vec2}, with no targeted SV evaluation. The Sri Lankan
deployment context further compounds the difficulty: bilingual and
code-switched speech is normal in everyday usage, telephony channels
dominate practical deployment, and labelled speakers are expensive to
collect.

\section{Problem statement}
\label{sec:problem}

This thesis addresses the question of how to construct a speaker
verification system that performs reliably on Sinhala, Tamil, and
code-switched audio when only a small labelled corpus is available.
We frame the problem along three dimensions:

\begin{enumerate}
    \item \textbf{Architectural choice.} Modern SV architectures
    (ECAPA-TDNN \cite{desplanques2020ecapa}, ResNet-SE variants,
    self-supervised front-ends) were validated on English VoxCeleb
    scale. Whether their inductive biases hold for low-resource
    Sinhala/Tamil is an open empirical question.
    \item \textbf{Cross-lingual transfer.} When a target-language
    corpus is small ($\leq$~10\,000 utterances), training from
    scratch overfits. Cross-lingual fine-tuning from an English-
    pretrained checkpoint is the natural alternative; the right
    freezing, learning-rate, and regularisation policy for SV
    fine-tuning specifically remains under-studied.
    \item \textbf{Calibration and operating point.} Per-language
    decision thresholds, score normalisation (AS-Norm
    \cite{matejka2017analysis}), generative scoring backends (PLDA),
    and language-aware regularisation (domain adversarial training
    \cite{ganin2015dann}) are standard tools for cross-lingual SV in
    the literature, but their interactions on Sinhala/Tamil data
    are not characterised.
\end{enumerate}

\section{Research questions}
\label{sec:rqs}

We address the following research questions:

\begin{description}
    \item[\textbf{RQ1}] How well do released, English/VoxCeleb-trained
    SV models transfer \emph{zero-shot} to Sinhala and Tamil, and how
    large is the degradation relative to their published English
    benchmarks?
    \item[\textbf{RQ2}] Which fine-tuning policy — frozen-encoder
    probing, parameter-efficient adaptation (LoRA), full fine-tuning,
    or layer-wise LR decay — best adapts a pretrained speech
    foundation model to Sinhala/Tamil per trainable parameter?
    \item[\textbf{RQ3}] Does language-adversarial training (gradient
    reversal / DANN \cite{ganin2015dann}) improve cross-lingual
    verification of bilingual speakers, and by how much?
    \item[\textbf{RQ4}] At eval time, do score normalisation (AS-Norm)
    and a generative scoring backend (PLDA) compose constructively, or
    do they redundantly model the same residual variance?
    \item[\textbf{RQ5}] How large is the gap between monolingual
    (Sinhala-only, Tamil-only) and cross-lingual EER for a single
    trained model, and can it be closed via threshold calibration
    alone or does it require architectural intervention?
\end{description}

\section{Contributions}
\label{sec:contributions}

The contributions of this thesis are:

\begin{enumerate}
    \item \textbf{A reproducible Sinhala/Tamil SV benchmark.}
    Benchmark v0 covers 527 speakers (478 Sinhala from OpenSLR SLR52,
    49 Tamil from OpenSLR SLR65) with seed-reproducible monolingual
    trial lists (12{,}000 trials per language, gender-controlled
    impostors where metadata permits), produced by a released
    deterministic pipeline; benchmark v1 extends the evaluation to
    SLCeleb \cite{slceleb2025}, our group's 280-speaker in-the-wild
    corpus with multi-session and bilingual speakers.
    \item \textbf{A comprehensive empirical study} of cross-lingual
    transfer for SV under low-resource constraints, in three suites:
    the first zero-shot baseline table for Sinhala (four modern
    architectures); a training-free adaptation ablation (AS-Norm,
    PLDA, per-language calibration) that quantifies how much of the
    cross-lingual degradation is a fixable score-domain problem; and
    a four-policy fine-tuning comparison (frozen / LoRA / full /
    layer-wise LR decay) of a WavLM foundation model, replicated
    over three seeds.
    \item \textbf{Per-language and pooled evaluation infrastructure}
    that supports principled comparison of Sinhala-only,
    Tamil-only, and cross-lingual trial sets in a single trainer
    invocation, with paper-grade statistical reporting
    (mean$\pm$std over three random seeds).
    \item \textbf{A practical open-source toolkit} extending the
    \texttt{voxceleb\_trainer} codebase \cite{slspv2026repo} with
    eleven new features (FEATURE-001..FEATURE-011) and twenty-six
    bug fixes (BUGFIX-001..BUGFIX-026). Each feature ships with a
    design document and is selectable via a single YAML key, enabling
    independent ablation.
\end{enumerate}

\section{Thesis structure}
\label{sec:structure}

The remainder of the thesis is organised as follows.
Chapter~\ref{ch:related} reviews prior work on speaker verification
architectures, cross-lingual transfer, score normalisation, and
language-aware modelling.
Chapter~\ref{ch:methods} describes the corpus, architectures, training
recipe, scoring backends, and evaluation protocol.
Chapter~\ref{ch:results} reports the headline and ablation results.
% Chapters 5 (Discussion) and 6 (Conclusion) are not scaffolded yet;
% switch the plain numbers below back to \ref{ch:discussion} /
% \ref{ch:conclusion} once those files exist.
Chapter~5 analyses the findings, addresses the
research questions of Section~\ref{sec:rqs}, and discusses limitations.
Chapter~6 summarises the contributions and outlines
future work, including the SLCeleb (v1) open-set benchmark, the
bilingual cross-lingual study, and self-supervised pretraining on
unlabelled Sri Lankan audio.

% ====================================================================
%  Chapter 2 — Related Work
% --------------------------------------------------------------------
%  Template chapter. The section headings cover the four sub-fields
%  the thesis touches:
%    2.1 Speaker verification architectures
%    2.2 Loss functions for speaker discrimination
%    2.3 Cross-lingual / low-resource transfer
%    2.4 Score normalisation and calibration
%    2.5 Multilingual self-supervised speech representation
%    2.6 Positioning of this work
%  Citations resolve against the .bib block in RUN_GUIDE.md §16.
%  Add references to .bib as you cite them; bracketed placeholders
%  like [CITE_DESPLANQUES_PAPER_PAGE] mark spots where a specific page
%  number or quoted figure should be added during the lit-review pass.
% ====================================================================

\chapter{Related Work}
\label{ch:related}

\section{Speaker verification architectures}
\label{sec:related-architectures}

The dominant architectures for modern speaker verification fall into
three families. The first is the residual convolutional family,
typified by ResNet-SE34V2 with squeeze-and-excitation blocks and
self-attentive pooling. The second is the time-delay neural network
(TDNN) family, dominated since 2020 by ECAPA-TDNN
\cite{desplanques2020ecapa}. ECAPA-TDNN introduced three key
ingredients: (i) Res2Net-style multi-scale convolutions with
hierarchical accumulation, (ii) squeeze-and-excite channel
re-weighting at each block, and (iii) multi-layer feature
aggregation followed by channel-attentive statistical pooling.
ECAPA-TDNN consistently reports 0.8--1.5\,\% EER on VoxCeleb1-O and
has become the de-facto reference architecture in subsequent SV
papers. The third family is self-supervised, replacing the
mel-front-end with a frozen pretrained transformer encoder such as
WavLM \cite{chen2022wavlm}, XLS-R, or mHuBERT-147; these models
typically outperform mel-based architectures on low-resource
language data by 30--50\,\% relative EER, at the cost of significant
parameter count (95M--300M).

% [CITE_OR_EXPAND: brief mention of x-vector / i-vector predecessors
% to ground the field's evolution; one paragraph is enough for an
% MPhil chapter.]

\section{Loss functions for speaker discrimination}
\label{sec:related-loss}

Speaker-verification training has progressively replaced cross-entropy
softmax with margin-based variants. AM-Softmax and Additive Angular
Margin Softmax (AAM-Softmax) \cite{deng2019arcface} reshape the
classification objective to produce more discriminative embeddings
on the hypersphere by adding a margin to the cosine similarity
between the embedding and the target-class weight vector. The
margin-based losses became standard in SV largely because they
produce embeddings that score well under cosine distance at eval
time — a closer match to how SV systems are actually used than the
intra-class compactness produced by plain softmax.

% [CITE_OR_EXPAND: mention sub-centre AAM-Softmax (\S 3.2 \#10 of the
% analysis doc) as a recent refinement; one paragraph if the SC-AAM
% ablation is included, omit otherwise.]

Auxiliary use of language information has also proven productive,
either during training via multi-task language-ID heads, or at
decision time via language-aware scoring and calibration
\cite{thienpondt2022calibration}.
Adversarial counterparts — Domain-Adversarial Neural Networks
\cite{ganin2015dann} attached to a language or channel classifier
through a gradient-reversal layer — pursue the opposite goal: forcing
the embedding to be \emph{uninformative} about the domain. Both
approaches appear in the SV literature; their relative effectiveness
depends on whether language is a useful signal (multi-task wins)
or a confound (adversarial wins).

\section{Cross-lingual and low-resource transfer}
\label{sec:related-crosslingual}

When a target-language SV corpus is small ($\leq$~10k utterances),
training a deep architecture from scratch overfits catastrophically.
The standard remedy is cross-lingual fine-tuning from a checkpoint
pretrained on a high-resource corpus (typically VoxCeleb1 or
VoxCeleb2). Three knobs control how aggressively the target
language overwrites the pretrained representation:

\begin{enumerate}
    \item \textbf{Selective freezing.} Freezing the front-end
    (SincConv or mel-filterbank parameters) preserves the
    acoustic-feature extractor learned on the source language. For
    very small corpora, freezing the entire encoder and training
    only the projection head is the safest option, and already
    captures most of the self-supervised representation's value
    \cite{chen2022wavlm,sang2025unipet}.
    \item \textbf{Learning-rate scaling.} A standard
    multiplier of 0.1$\times$ to 0.01$\times$ the from-scratch LR
    avoids destabilising the prior.
    \item \textbf{Layer-wise learning-rate decay (LLRD).} A decay
    factor $\gamma \in [0.8, 0.95]$ applied per layer from the
    encoder top to the input front-end, so that deeper layers carry
    transferable acoustic priors while shallower (closer-to-output)
    layers adapt more aggressively. LLRD is universal in the BERT
    and WavLM fine-tuning literature; its application to SV
    fine-tuning is less systematically reported, motivating one of
    this thesis's ablations.
\end{enumerate}

% [CITE_OR_EXPAND: position prior cross-lingual SV papers — e.g. NIST
% SRE multilingual evaluations, MLS, MUSAN-based low-resource SV
% work — relative to the SL setting.]

\section{Score normalisation and calibration}
\label{sec:related-scoring}

Two eval-time mechanisms address the gap between training-time and
deployment-time score distributions. Adaptive Symmetric Score
Normalisation (AS-Norm) \cite{matejka2017analysis} computes
per-trial-side cohort statistics from a fixed set of impostor
embeddings and normalises trial scores against those statistics.
It is universal in the NIST SRE pipeline and typically yields
5--20\,\% relative MinDCF reduction at zero training cost.

Probabilistic Linear Discriminant Analysis (PLDA) \cite{kenny2010,
sizov2014unifying} models the embedding distribution generatively
and replaces cosine scoring with a likelihood ratio test under
``same speaker'' vs ``different speaker'' hypotheses. PLDA was the
dominant SV backend in the i-vector era and remains useful for
short-utterance and channel-mismatched trials, although on modern
discriminatively-trained embeddings the marginal gain over cosine is
smaller. The two-covariance simplification of PLDA
\cite{sizov2014unifying} provides nearly identical EER to the full
formulation at half the computational complexity.

Score normalisation and PLDA compose: AS-Norm normalises against
cohort statistics regardless of whether the underlying score is
cosine or PLDA. The NIST SRE pipeline order is PLDA $\to$ AS-Norm.

% [CITE_OR_EXPAND: per-language threshold calibration is a separate
% research direction (Platt scaling, isotonic regression). One short
% paragraph for completeness.]

\section{Multilingual self-supervised speech representation}
\label{sec:related-ssl}

Self-supervised models trained on large multilingual audio corpora
(wav2vec 2.0 \cite{baevski2020wav2vec2}, WavLM
\cite{chen2022wavlm}, XLS-R, mHuBERT-147 \cite{boito2024mhubert})
provide speech-aware embeddings without language-specific labelling.
mHuBERT-147 is particularly relevant to the Sri Lankan setting: it
deliberately upweights low-resource languages including Sinhala and
Tamil in its pretraining corpus.

There are three ways to use such models for low-resource SV:
(i) load an off-the-shelf checkpoint, freeze it, and train only a
pooling head and projection on the target-language labelled data
% [CITE_OR_EXPAND: typical SV-on-WavLM results]; (ii) continue
pretraining (CPT) on unlabelled target-language audio before
SV-fine-tuning, which adds 5--15\,\% relative EER on top of (i)
when sufficient unlabelled audio ($\geq$~500\,h) is available
\cite{pratap2024mms}; (iii) pretrain from scratch on
target-language-only audio, which requires thousands of hours and
rarely outperforms (ii). This thesis evaluates option (i) as the
SSL baseline; option (ii) is documented as future work.

\section{Positioning of this work}
\label{sec:positioning}

The contributions outlined in
Section~\ref{sec:contributions} sit at the intersection of three
otherwise separately-developed lines of work:

\begin{itemize}
    \item \textbf{Architectures and losses} for English-scale SV
    (\S~\ref{sec:related-architectures}, \S~\ref{sec:related-loss})
    are applied to low-resource Sinhala/Tamil with no targeted
    benchmark in the literature.
    \item \textbf{Cross-lingual transfer techniques}
    (\S~\ref{sec:related-crosslingual}) are widely studied for ASR
    and translation but less so for SV; their interaction with
    SV-specific losses (AAM-Softmax) and SV-specific eval backends
    (PLDA, AS-Norm) has not been systematically ablated.
    \item \textbf{Multilingual SSL encoders}
    (\S~\ref{sec:related-ssl}) cover Sinhala and Tamil only as part
    of broader low-resource batches, with no published SV evaluation
    on either language.
\end{itemize}

To the best of our knowledge, this is the first study to provide a
reproducible SV benchmark on Sinhala and Tamil with full ablation
isolation of the cross-lingual transfer levers. The benchmark, the
corpus preparation pipeline, and the open-source toolkit
(\S~\ref{sec:contributions}) are released to support future work in
this direction.

% ====================================================================
%  Chapter 3 — Methods
% --------------------------------------------------------------------
%  Rewritten 2026-07-03 to describe the experiments actually run
%  (benchmark v0 + zero-shot baselines + training-free backend
%  adaptation + WavLM PEFT). Sources of record:
%    research_logs/2026-07-03-sl-benchmark-v0-design.md
%    research_logs/2026-07-03-zeroshot-baseline-table.md
%    research_logs/2026-07-03-p2-backend-adaptation.md
%    research_logs/2026-07-03-p3-peft-finetuning.md
%  The v1 (SLCeleb) protocol will be added once the corpus is staged.
% ====================================================================

\chapter{Methods}
\label{ch:methods}

\section{Corpus and benchmark construction}
\label{ssec:corpus}

We evaluate on a two-tier benchmark. \textbf{Benchmark v0}, used for
all experiments in this thesis, is built from two public read-speech
corpora: OpenSLR SLR52 \cite{kjartansson2018slr52} contributes 478
Sinhala speakers ($\approx$185k crowdsourced utterances, 16\,kHz FLAC,
minimum 10 utterances per speaker), and OpenSLR SLR65
\cite{he2020opensource} contributes 49 Tamil speakers (4{,}286 studio
utterances, 48\,kHz resampled at load). Speaker identities come from
the corpora's published metadata; gender labels exist only for SLR65.
A deterministic ingestion script (\texttt{tools/ingest\_openslr.py})
lays the corpus out as
\texttt{<lang>/<speaker>/<utterance>} symlinks and emits a manifest
and speaker-metadata table. \textbf{Benchmark v1}, pending data
staging, replaces the evaluation side with SLCeleb
\cite{slceleb2025}, our group's 280-speaker in-the-wild
Sinhala/Tamil celebrity corpus, which adds multi-genre audio,
multi-session speakers, and bilingual speakers for true
cross-lingual trials.

\section{Trial construction}
\label{ssec:trials}

Utterances are split 80/20 per speaker into train and test partitions
(seed 42, \texttt{tools/sl\_dataprep.py}). From the test partition we
sample, per language, 3{,}000 same-speaker target pairs and 9{,}000
same-language impostor pairs (1:3 ratio), with unordered-pair
deduplication and attempt-bounded sampling. Impostor pairs are
same-gender whenever both speakers' genders are known; this holds for
all Tamil trials but cannot be enforced for Sinhala (SLR52 has no
gender metadata). Cross-session targets are enforced when session
keys are derivable; the OpenSLR layouts carry no session metadata, so
v0 targets are not session-controlled. Consequently v0 absolute error
rates are optimistic and are valid for \emph{relative} comparison
between systems, not for comparison against in-the-wild benchmarks;
we state this caveat wherever v0 numbers are reported. The
cross-lingual trial list is empty at v0 (the OpenSLR corpora contain
no bilingual speakers) and becomes meaningful at v1.

\section{Metrics}
\label{ssec:metrics}

We report equal error rate computed as the mean of the false-positive
and false-negative rates at the crossing point (EER$_{\mathrm{avg}}$),
and minimum detection cost (minDCF) at two operating points,
$p_{\mathrm{target}} \in \{0.01, 0.05\}$ with
$C_{\mathrm{miss}}=C_{\mathrm{fa}}=1$, following NIST and VoxCeleb
practice. Calibration quality is measured with $C_{\mathrm{llr}}$ and
actual DCF after logistic calibration. Unlabeled trials are treated
as a hard error by the evaluation harness.

\section{Experiment suite A: zero-shot baselines}
\label{ssec:zeroshot}

Four released, VoxCeleb-trained checkpoints are evaluated with no
adaptation of any kind: ECAPA-TDNN \cite{desplanques2020ecapa} as
distributed by SpeechBrain \cite{ravanelli2021speechbrain}
($\approx$20M parameters), and ReDimNet-B1/B2/B6
\cite{yakovlev2024redimnet} (2.2M/4.7M/15M). Scoring is cosine
similarity on L2-normalised full-utterance embeddings, with no score
normalisation and no calibration; this table is the raw-transfer
reference that all adaptation levers are measured against.

\section{Experiment suite B: training-free backend adaptation}
\label{ssec:backend}

On top of cached suite-A embeddings we ablate three eval-time levers,
none of which updates the encoder:
(i)~\textbf{AS-Norm} \cite{matejka2017analysis}: adaptive symmetric
score normalisation with a 527-utterance in-language cohort (one
utterance per training-pool speaker), top-$K{=}300$;
(ii)~\textbf{PLDA} \cite{sizov2014unifying,kenny2010}: a
two-covariance model ($d{=}150$) trained on 10{,}480 utterances
($\leq$20 per speaker) from the training partition, preceded by
centering, LDA, and length normalisation;
(iii)~\textbf{language-aware logistic calibration}: an affine
score-to-LLR map fitted per language on half the trials and evaluated
on the other half, plus the cross-language transplant of each
calibrator to quantify the cross-lingual score shift
\cite{thienpondt2022calibration}. The grid crosses
$\{$cosine, PLDA$\}\times\{$raw, AS-Norm$\}$ for three encoders,
answering RQ4 directly. The PLDA fit necessarily sees the trial
speakers (utterance-disjoint, standard for domain-adaptation PLDA);
we flag this closed-set property explicitly.

\section{Experiment suite C: parameter-efficient fine-tuning}
\label{ssec:peft}

Suite C fine-tunes WavLM-base+ \cite{chen2022wavlm} (94.4M
parameters, English-pretrained) for Sinhala/Tamil speaker
verification. The 13 hidden states are combined by a learnable
softmax-weighted sum, pooled by attentive statistics pooling (128-d
bottleneck), and projected to a 192-d embedding trained with
AAM-Softmax \cite{deng2019arcface} (margin 0.2, scale 30) over the
527 training speakers. The CNN feature extractor stays frozen in all
arms. Four fine-tuning policies are compared:

\begin{description}
  \item[frozen] encoder frozen; only the pooling/embedding head and
  layer weights train ($\approx$0.5\% of parameters);
  \item[LoRA] rank-8 adapters ($\alpha{=}16$) \cite{hu2022lora} on the
  query/key/value/output projections of every attention block, plus
  the head (1.08M trainable, 1.1\%);
  \item[full] all encoder parameters unfrozen (encoder LR $10^{-5}$,
  head LR $10^{-3}$);
  \item[LLRD] full fine-tuning with layer-wise learning-rate decay
  ($\gamma{=}0.9$ per layer from the top).
\end{description}

Training uses the v0 train split capped at 60 utterances per speaker
(31{,}049 utterances), random 2.0\,s crops, batch 64, AdamW (weight
decay $2{\times}10^{-5}$), exponential LR decay ($\gamma{=}0.95$),
mixed precision, gradient clipping at 5.0, and 15 epochs. Model
selection uses per-epoch pooled EER on stratified 2{,}000-pair
subsamples; the selected checkpoint is finally evaluated on the full
12{,}000-pair lists, making suite-C numbers directly comparable to
suites A and B. Because the evaluation speakers are the training
speakers (utterance-disjoint split), suite C measures
\emph{representation adaptation} rather than open-set
generalisation; the v1 rerun provides the open-set condition.
Finally, the suite-B backend stack is re-applied to the fine-tuned
embeddings (suite B$\times$C composition) to test whether backend
adaptation retains value once the encoder is in-domain
\cite{wang2022plda}.

\section{Reproducibility}
\label{ssec:repro}

Trial lists, cohort lists, and PLDA training lists are reproducible
from seed 42. The LoRA and full arms of suite C are replicated with
three seeds (42/43/44) and reported as mean$\pm$std; the frozen and
LLRD arms use seed 42 only, as both sit far from the decision
boundary between policies. All runs, configurations, per-epoch
histories, and raw result JSONs are archived with the research logs
listed in the chapter header.

% ====================================================================
%  Chapter 4 — Results
% --------------------------------------------------------------------
%  Rewritten 2026-07-03 with the completed benchmark-v0 results.
%  Every number is sourced from:
%    research_logs/2026-07-03-zeroshot-baseline-table.md   (Table A)
%    research_logs/2026-07-03-p2-backend-adaptation.md     (Tables B, C, E)
%    research_logs/2026-07-03-p3-peft-finetuning.md        (Table D)
%  Raw JSONs: /home/anuraj/sl_spv_bench/results/ and runs/p3_*/.
%  All numbers carry the v0 caveat of Section 3.2 (read speech,
%  session-uncontrolled targets, closed-set): valid for relative
%  comparison; v1 (SLCeleb) supplies the open-set condition.
% ====================================================================

\chapter{Results}
\label{ch:results}

\section{Zero-shot baselines (suite A)}
\label{sec:zeroshot-results}

Table~\ref{tab:zeroshot} reports, to our knowledge, the first
published speaker-verification results for Sinhala. All four
VoxCeleb-trained checkpoints transfer at usable but clearly degraded
accuracy: relative to each model's published VoxCeleb1-O EER, Sinhala
EER inflates by $4.9$--$8.0\times$ and Tamil by $2.3$--$4.0\times$ —
Tamil within the $2$--$4\times$ band the CN-Celeb cross-domain
literature reports, Sinhala well above it, even though v0 is clean
read speech. The VoxCeleb ranking is mostly
preserved (ReDimNet-B6 best, ECAPA worst) with one inversion:
ReDimNet-B1 beats the nominally stronger B2 on both languages, and
with 2.2M parameters also beats the $\approx$20M-parameter ECAPA —
parameter count does not monotonically buy language robustness.
Across all systems minDCF($10^{-2}$) is poor (0.24--0.48) while EERs
look usable: the signature of badly scaled scores in a mismatched
domain, which motivates suite B.

\begin{table}[t]
\centering
\caption{Zero-shot transfer of VoxCeleb-trained checkpoints to
benchmark v0. EER$_{\mathrm{avg}}$\,\% and minDCF at
$p_{\mathrm{t}}{=}0.01/0.05$; 12{,}000 trials per language.
$\times$Vox1-O = SL EER divided by the checkpoint's published
VoxCeleb1-O EER.}
\label{tab:zeroshot}
\begin{tabular}{lcccccc}
\toprule
& \multicolumn{3}{c}{\textbf{Sinhala}} & \multicolumn{3}{c}{\textbf{Tamil}} \\
\cmidrule(lr){2-4}\cmidrule(lr){5-7}
\textbf{Model} & EER & minDCF & $\times$Vox1-O & EER & minDCF & $\times$Vox1-O \\
\midrule
ECAPA-TDNN (SpeechBrain) & 4.78 & 0.471/0.291 & 6.0 & 3.21 & 0.482/0.207 & 4.0 \\
ReDimNet-B1              & 3.57 & 0.306/0.209 & 4.9 & 1.67 & 0.406/0.162 & 2.3 \\
ReDimNet-B2              & 4.17 & 0.422/0.244 & 8.0 & 1.97 & 0.454/0.170 & 3.8 \\
ReDimNet-B6              & \textbf{2.71} & 0.238/0.153 & 7.3 & \textbf{1.48} & 0.422/0.160 & 4.0 \\
\bottomrule
\end{tabular}
\end{table}

\section{Training-free backend adaptation (suite B)}
\label{sec:backend-results}

Table~\ref{tab:backend} crosses the scoring backend (cosine vs PLDA)
with AS-Norm for three encoders. Three findings:

\begin{enumerate}
\item \textbf{AS-Norm delivers 12--39\% relative EER reduction}
(Sinhala: 33/33/23\% for ECAPA/B1/B6; Tamil: 30/12/39\%), bracketing
the 20--30\% band reported by \cite{matejka2017analysis}, and
improves minDCF($10^{-2}$) in five of the six model--language cells
(best: B6 Sinhala $0.238\!\rightarrow\!0.167$; the exception is B1
Tamil, $0.406\!\rightarrow\!0.433$). It is the most cost-effective
lever in this study: no training, one cohort file.
\item \textbf{PLDA hurts Sinhala but helps Tamil.} PLDA increases
Sinhala EER for every encoder (B6: $2.71\!\rightarrow\!3.84$) while
reducing Tamil EER for two of the three encoders (ECAPA:
$3.21\!\rightarrow\!2.46$; B6: $1.48\!\rightarrow\!1.11$; B1 worsens
slightly, $1.67\!\rightarrow\!1.77$). Our working hypothesis is a benchmark
artifact: the PLDA fit is dominated by the 478 Sinhala speakers of
near-single-session crowdsourced audio, so its within-class
covariance underestimates true channel variability, whereas the
studio-recorded Tamil fits the model. The v1 multi-session rerun is
the decisive test before dismissing the ``PLDA wins under shift''
claim of \cite{wang2022plda}.
\item \textbf{RQ4 (composition):} at v0, AS-Norm dominates PLDA and
the composed stack does not beat cosine{+}AS-Norm on Sinhala, though
it yields the best Tamil minDCF. The answer is
condition-dependent; both levers stay in the v1 protocol.
\end{enumerate}

The best training-free v0 system is
\textbf{ReDimNet-B6 + AS-Norm: 2.07\% (si), 0.90\% (ta)}.

\begin{table}[t]
\centering
\caption{Suite B ablation: EER\,\% (minDCF at $10^{-2}$).
Cosine rows reproduce Table~\ref{tab:zeroshot} exactly (sanity
check). Bold = best per encoder and language.}
\label{tab:backend}
\begin{tabular}{llcc}
\toprule
\textbf{Encoder} & \textbf{Backend} & \textbf{Sinhala} & \textbf{Tamil} \\
\midrule
ECAPA-TDNN & cosine          & 4.78 (0.471) & 3.21 (0.482) \\
           & cosine+AS-Norm  & \textbf{3.19} (0.330) & 2.24 (0.480) \\
           & PLDA            & 6.00 (0.598) & 2.46 (0.495) \\
           & PLDA+AS-Norm    & 5.80 (0.482) & \textbf{2.06} (0.374) \\
\midrule
ReDimNet-B1 & cosine         & 3.57 (0.306) & 1.67 (0.406) \\
            & cosine+AS-Norm & \textbf{2.38} (0.245) & 1.47 (0.433) \\
            & PLDA           & 5.34 (0.452) & 1.77 (0.419) \\
            & PLDA+AS-Norm   & 5.23 (0.382) & \textbf{1.49} (0.297) \\
\midrule
ReDimNet-B6 & cosine         & 2.71 (0.238) & 1.48 (0.422) \\
            & cosine+AS-Norm & \textbf{2.07} (0.167) & \textbf{0.90} (0.376) \\
            & PLDA           & 3.84 (0.293) & 1.11 (0.340) \\
            & PLDA+AS-Norm   & 3.56 (0.254) & 1.16 (0.314) \\
\bottomrule
\end{tabular}
\end{table}

\subsection{Cross-lingual score shift and calibration}
\label{sec:calibration-results}

Table~\ref{tab:crosscal} fits an affine logistic calibrator per
language and then transplants it across languages (ReDimNet-B6).
Self-calibration achieves good $C_{\mathrm{llr}}$ (0.066--0.171), but
the cross-language transplant inflates $C_{\mathrm{llr}}$ by up to
$5.7\times$ (Sinhala scored with a Tamil-fit calibrator, cosine
{+}AS-Norm). The cross-lingual score shift is therefore large and
real, extending \cite{thienpondt2022calibration} to Sinhala/Tamil,
and \textbf{per-language calibration is mandatory for deployment}.
A publishable nuance: PLDA scores are the most \emph{transportable}
across languages (penalty only $1.2$--$2.7\times$) even where PLDA's
EER is worse — the generative backend produces more
language-homogeneous score scales than cosine.

\begin{table}[t]
\centering
\caption{Language-aware calibration, ReDimNet-B6:
$C_{\mathrm{llr}}$ with a self-language-fit calibrator vs a
calibrator transplanted from the other language.}
\label{tab:crosscal}
\begin{tabular}{lcccc}
\toprule
& \multicolumn{2}{c}{\textbf{Sinhala}} & \multicolumn{2}{c}{\textbf{Tamil}} \\
\cmidrule(lr){2-3}\cmidrule(lr){4-5}
\textbf{Backend} & self & ta-fit & self & si-fit \\
\midrule
cosine          & 0.129 & 0.489 ($3.8\times$) & 0.095 & 0.149 ($1.6\times$) \\
cosine+AS-Norm  & 0.099 & 0.565 ($5.7\times$) & 0.071 & 0.167 ($2.4\times$) \\
PLDA            & 0.171 & 0.390 ($2.3\times$) & 0.073 & 0.084 ($1.2\times$) \\
PLDA+AS-Norm    & 0.163 & 0.442 ($2.7\times$) & 0.066 & 0.094 ($1.4\times$) \\
\bottomrule
\end{tabular}
\end{table}

\section{Parameter-efficient fine-tuning (suite C)}
\label{sec:peft-results}

Table~\ref{tab:peft} compares the four fine-tuning policies.
Accuracy is monotone in encoder freedom: frozen $\prec$ LoRA
$\prec$ LLRD $\prec$ full. Full fine-tuning of WavLM-base+ on 31k
in-language utterances is the best system in this thesis —
\textbf{1.69$\pm$0.06\% (si) / 1.41$\pm$0.10\% (ta)} over three
seeds — beating the best zero-shot supervised model by 38\% relative
on Sinhala (2.71\% $\rightarrow$ 1.69\%) and matching it within noise
on Tamil.

Two observations qualify the PEFT narrative:

\begin{enumerate}
\item \textbf{The ``PEFT $\geq$ full fine-tuning under shift''
prediction \cite{sang2025unipet} does not reproduce here}: full
fine-tuning beats LoRA at every seed, with a $2.2\times$ mean EER
gap (si 1.69 vs 3.73). A plausible explanation is the closed-set v0
protocol (Section~\ref{ssec:peft}): full fine-tuning's usual
penalty — overfitting the adaptation speakers at the cost of unseen
speakers — is not measured when the evaluation speakers \emph{are}
the training speakers. The open-set v1 rerun arbitrates. Secondary
factors: LoRA $r{=}8$ may be undersized and its head LR untuned.
\item \textbf{LoRA is high-variance across seeds}
($\pm$0.49--0.61 EER, vs $\pm$0.06--0.10 for full): single-seed
PEFT-vs-full comparisons on corpora of this size are unsound, and
the seed-42-only column would have flattered LoRA with its best
draw.
\end{enumerate}

Even the frozen probe (0.5\% of parameters) lands within
$\approx$1 EER point of zero-shot ReDimNet-B6 on Sinhala — the
pretrained WavLM features already carry most of the speaker signal,
and everything beyond frozen is adaptation gain.

\begin{table}[t]
\centering
\caption{Suite C: fine-tuning policy comparison on benchmark v0
(EER$_{\mathrm{avg}}$\,\%, full 12{,}000-pair lists, cosine scoring).
Seed-42 column for all arms; mean$\pm$std over seeds 42/43/44 for
LoRA and full. Zero-shot reference: ReDimNet-B6 2.71 (si) / 1.48 (ta).}
\label{tab:peft}
\begin{tabular}{lccccc}
\toprule
& & \multicolumn{2}{c}{\textbf{seed 42}} & \multicolumn{2}{c}{\textbf{mean$\pm$std (3 seeds)}} \\
\cmidrule(lr){3-4}\cmidrule(lr){5-6}
\textbf{Arm} & \textbf{Trainable} & si & ta & si & ta \\
\midrule
frozen & 0.5\%  & 3.74 & 3.13 & --- & --- \\
LoRA   & 1.1\%  & 3.17 & 2.44 & 3.73$\pm$0.49 & 3.11$\pm$0.61 \\
LLRD   & 95.6\% & 2.40 & 1.81 & --- & --- \\
full   & 95.6\% & \textbf{1.63} & \textbf{1.30} & \textbf{1.69$\pm$0.06} & \textbf{1.41$\pm$0.10} \\
\bottomrule
\end{tabular}
\end{table}

\section{Composing backend adaptation with fine-tuning
(suite B$\times$C)}
\label{sec:composition-results}

Re-applying the suite-B stack to the fine-tuned (full, seed 42)
embeddings tests whether backend adaptation survives encoder
adaptation (Table~\ref{tab:composition}). AS-Norm's EER gain
saturates (si $-0.07$ absolute; ta still $+13\%$ relative) but it
keeps improving the operating-point metrics — minDCF($10^{-2}$)
si $0.180\!\rightarrow\!0.144$, ta $0.175\!\rightarrow\!0.149$, and
$C_{\mathrm{llr}}$ down $\approx$15\% on both — so it remains worth
deploying even with an in-domain encoder. PLDA, by contrast, is now
clearly redundant-to-harmful on both languages, exactly the
in-domain prediction of \cite{wang2022plda}: fine-tuning moved the
embedding distribution in-domain, and the v0 PLDA story is coherent
before and after. The best overall v0 operating system is
\textbf{fine-tuned WavLM (full) + AS-Norm}: si 1.71\%/0.144, ta
1.13\%/0.149 (or plain cosine when EER is the sole criterion).

\begin{table}[t]
\centering
\caption{Suite B backends applied to fine-tuned WavLM embeddings
(full arm, seed 42): EER\,\% (minDCF at $10^{-2}$).}
\label{tab:composition}
\begin{tabular}{lcc}
\toprule
\textbf{Backend} & \textbf{Sinhala} & \textbf{Tamil} \\
\midrule
cosine          & \textbf{1.63} (0.180) & 1.30 (0.175) \\
cosine+AS-Norm  & 1.71 (\textbf{0.144}) & \textbf{1.13} (\textbf{0.149}) \\
PLDA            & 2.46 (0.258) & 1.57 (0.234) \\
PLDA+AS-Norm    & 2.37 (0.197) & 1.43 (0.201) \\
\bottomrule
\end{tabular}
\end{table}

\section{Per-language analysis}
\label{sec:per-lang}

Tamil EER is uniformly lower than Sinhala (studio recordings, 49
speakers, gender-controlled impostors), yet Tamil
minDCF($10^{-2}$) is uniformly \emph{worse} for the ReDimNet family
(e.g.\ B6: 0.238 si vs 0.422 ta). With only 49 speakers the impostor
pool is small, and a few hard same-genre impostors dominate the
low-false-alarm operating region. Two methodological consequences:
EER alone overstates how ``solved'' the Tamil condition is, and both
metrics must be read against the corpus asymmetry (478 crowdsourced
phone-quality Sinhala speakers vs 49 studio Tamil speakers) rather
than interpreted as a property of the languages themselves.
Cross-lingual (enrol-in-one-language, test-in-the-other) trials
require bilingual speakers and therefore await benchmark v1.

\section{Statistical reporting}
\label{sec:significance}

Multi-seed replication (42/43/44) covers the two arms whose ordering
constitutes the headline claim (LoRA vs full); their separation
(at least 1.1 EER points at every matched seed, and non-overlapping
seed ranges on both languages) makes the verdict robust without
interval estimation. The frozen and LLRD
arms, and all training-free results, are deterministic given the
seed-42 lists. For the journal version, bootstrap confidence
intervals over trial pairs (1{,}000 resamples, 2.5/97.5 quantiles)
remain planned once the v1 benchmark is in place.


% ---- Chapter 5 — Discussion (not yet scaffolded) ----
% \include{thesis_chapters/05_discussion}

% ---- Chapter 6 — Conclusion (not yet scaffolded) ----
% \include{thesis_chapters/06_conclusion}

\bibliographystyle{plain}
\bibliography{thesis_chapters/references}

\end{document}
